% ============================================================
%  PREAMBLE -- Proteksi Sistem Tenaga Listrik (disalin dari TILI)
% ============================================================
\usepackage[a4paper,margin=2.2cm,top=2.3cm,bottom=2.5cm]{geometry}
\usepackage[T1]{fontenc}
\usepackage[utf8]{inputenc}
\usepackage{lmodern}
\usepackage[indonesian]{babel}
\usepackage{amsmath,amssymb}
\usepackage{icomma}
\usepackage{xcolor}
\usepackage[most]{tcolorbox}
\usepackage[american]{circuitikz}
\usepackage{pgfplots}
\pgfplotsset{compat=1.18}
\usepgfplotslibrary{fillbetween}
\usetikzlibrary{arrows.meta,calc,positioning,decorations.pathreplacing,patterns,angles,quotes,shapes.geometric}
\usepackage{booktabs,array,multirow,longtable,tabularx}
\usepackage{enumitem}
\usepackage{etoolbox}
\usepackage{fontawesome5}
\usepackage{caption}
\captionsetup{font=small,labelfont=bf}
\usepackage{float}
\usepackage[hidelinks]{hyperref}
\setlist{itemsep=1pt,topsep=3pt}
\setlength{\parindent}{0pt}
\setlength{\parskip}{4pt}
\renewcommand{\arraystretch}{1.18}

% ------------------------------------------------------------
%  PALET: 5 keluarga warna x 4 gradasi.
%  Makna tiap warna DITETAPKAN ULANG di awal setiap bab.
% ------------------------------------------------------------
\definecolor{biru1}{HTML}{0D47A1}  % biru tua
\definecolor{biru2}{HTML}{1E88E5}  % biru terang
\definecolor{biru3}{HTML}{00838F}  % biru kehijauan (sian)
\definecolor{biru4}{HTML}{3949AB}  % biru nila
\definecolor{merah1}{HTML}{B71C1C} % merah tua
\definecolor{merah2}{HTML}{E53935} % merah terang
\definecolor{merah3}{HTML}{EF6C00} % jingga
\definecolor{merah4}{HTML}{AD1457} % merah keunguan
\definecolor{hijau1}{HTML}{1B5E20} % hijau tua
\definecolor{hijau2}{HTML}{43A047} % hijau terang
\definecolor{hijau3}{HTML}{827717} % hijau zaitun
\definecolor{hijau4}{HTML}{00796B} % hijau toska
\definecolor{ungu1}{HTML}{4A148C}  % ungu tua
\definecolor{ungu2}{HTML}{8E24AA}  % ungu terang
\definecolor{ungu3}{HTML}{5E35B1}  % ungu kebiruan
\definecolor{ungu4}{HTML}{C2185B}  % ungu kemerahan (magenta)
\definecolor{coklat1}{HTML}{5D4037} % coklat tua
\definecolor{coklat2}{HTML}{A1572A} % coklat bata
\definecolor{coklat3}{HTML}{B8860B} % kuning tua / emas
\definecolor{coklat4}{HTML}{8D6E63} % coklat muda

% ------------------------------------------------------------
%  Simbol berwarna
%  \defq{nama}{warna}{simbol}{nama besaran}
%    -> \nama   : simbol berwarna (bisa di teks maupun rumus)
%    -> \tnama  : nama besaran, tebal & berwarna sama
% ------------------------------------------------------------
\newcommand{\defsym}[3]{\expandafter\gdef\csname #1\endcsname{\ensuremath{\textcolor{#2}{#3}}}}
\newcommand{\kw}[2]{\textcolor{#1}{\textbf{#2}}}
\newcommand{\defq}[4]{\defsym{#1}{#2}{#3}\expandafter\gdef\csname t#1\endcsname{\kw{#2}{#4}}}
\newcommand{\cl}[2]{\textcolor{#1}{#2}}   % warnai potongan rumus apa saja

% ------------------------------------------------------------
%  Kotak-kotak
% ------------------------------------------------------------
% Rumus kunci: bingkai tipis
\newtcbox{\kotak}{enhanced,colback=white,colframe=black!60,boxrule=0.5pt,arc=0pt,
  left=6pt,right=6pt,top=3pt,bottom=3pt,nobeforeafter,tcbox raise base}
% Setiap kotak rumus bernomor (bab.urutan), misalnya (2.5). Pasang \label{r:...} tepat sesudahnya.
\newcounter{rumus}[section]
\renewcommand{\therumus}{\thesection.\arabic{rumus}}
\newcommand{\nomorwarna}{biru1}
\newsavebox{\rumusbox}
\newcommand{\rumus}[2][]{\par\refstepcounter{rumus}\sbox{\rumusbox}{\kotak{%
  {\sffamily\bfseries\color{\nomorwarna}(\therumus)}\quad
  \if\relax\detokenize{#1}\relax\else{\scriptsize\sffamily\color{black!55}#1}\quad\fi
  $\displaystyle #2$}}\begin{center}\ifdim\wd\rumusbox>\linewidth
  \resizebox{\linewidth}{!}{\usebox{\rumusbox}}\else\usebox{\rumusbox}\fi\end{center}\par}

% Pemahaman penting: warna mengikuti keluarga konsepnya
\newtcolorbox{paham}[2][biru1]{enhanced,breakable,colback=#1!5,colframe=#1,
  boxrule=0pt,leftrule=3pt,arc=0pt,outer arc=0pt,
  colbacktitle=#1!13,coltitle=#1!80!black,fonttitle=\bfseries,
  title={\faLightbulb[regular]\enspace Pemahaman penting: #2},
  left=7pt,right=7pt,top=4pt,bottom=4pt,before skip=8pt,after skip=8pt}

% Proses di balik layar (penurunan / alasan): netral abu-abu
\newtcolorbox{balik}[1]{enhanced,breakable,colback=black!2.5,colframe=black!45,
  boxrule=0pt,leftrule=3pt,arc=0pt,outer arc=0pt,
  colbacktitle=black!8,coltitle=black!80,fonttitle=\bfseries,
  title={\faCogs\enspace Proses di balik layar: #1},
  left=7pt,right=7pt,top=4pt,bottom=4pt,before skip=8pt,after skip=8pt}

% Koreksi / catatan tentang slide: garis putus-putus
\newtcolorbox{koreksi}[1][Catatan tentang slide]{enhanced,breakable,colback=white,
  frame hidden,borderline={0.7pt}{0pt}{black!65,dashed},arc=0pt,
  fonttitle=\bfseries,coltitle=black,colbacktitle=white,
  title={\faExclamationTriangle\enspace #1},
  left=7pt,right=7pt,top=3pt,bottom=4pt,before skip=8pt,after skip=8pt}

% Contoh soal
\newtcolorbox{contoh}[1]{enhanced,breakable,colback=white,colframe=black!40,
  boxrule=0.5pt,arc=1pt,colbacktitle=black!8,coltitle=black,fonttitle=\bfseries,
  title={\faEdit[regular]\enspace Contoh Soal: #1},
  left=7pt,right=7pt,top=4pt,bottom=4pt,before skip=8pt,after skip=8pt}
\newcommand{\jawab}{\par\smallskip\textbf{Penyelesaian.}\ }

% Prosedur penyelesaian: Diketahui -> Ditanya -> Rencana -> Langkah 1, 2, ...
% Setiap langkah: \cari{besaran} lalu alasan (kenapa sekarang), rumus dalam simbol, baru angka.
\newcommand{\prosedurwarna}{biru3!75!black}
% Diketahui: butir dipisah titik koma (;), tiap butir satu baris.
\DeclareListParser*{\dkparse}{;}
\newcommand{\dkitem}[1]{\par\hangindent=2.2em\noindent\makebox[2.2em][l]{\hspace{0.9em}\textcolor{\prosedurwarna}{\textbullet}}\ignorespaces#1}
\newcommand{\diket}[1]{\par\smallskip\noindent\textbf{\textcolor{\prosedurwarna}{Diketahui.}}\par
  {\raggedright\dkparse{\dkitem}{#1}\par}}
\newcommand{\ditanya}[1]{\par\noindent\textbf{\textcolor{\prosedurwarna}{Ditanya.}}\ #1\par}
\newcommand{\rencana}[1]{\par\noindent\textbf{\textcolor{\prosedurwarna}{Rencana.}}\ #1\par}
\newcommand{\bagsoal}[1]{\par\smallskip\noindent\textbf{#1}\par}
% Langkah: judul "Langkah n. Cari X." lalu baris-baris berlabel di bawahnya.
\newcounter{lkh}
\newenvironment{langkah}{\par\setcounter{lkh}{0}\raggedright
  \setlength{\lineskiplimit}{3pt}\setlength{\lineskip}{4pt}
  \def\item{\par\addvspace{5pt}\stepcounter{lkh}\noindent
    {\bfseries\color{\prosedurwarna}Langkah \arabic{lkh}.}\ }}{\par\smallskip}
\newcommand{\cari}[1]{\textbf{Cari #1.}\ }
\newcommand{\cek}{\textbf{Cek.}\ }
\newlength{\lblw}\setlength{\lblw}{5.4em}
\newcommand{\baris}[2]{\par\hangindent=\lblw\hangafter=1\noindent
  \makebox[\lblw][l]{\hspace{1.3em}{\small\sffamily\color{black!55}#1}}#2\par}
\newcommand{\alasan}[1]{\baris{Kenapa}{#1}}
% rumus dari kotak bernomor: nomor + rumus ditulis ulang
\newcommand{\pakai}[2]{\baris{Rumus}{{\sffamily\bfseries\color{\nomorwarna}(\ref{#1})}\enspace$\displaystyle #2$}}
% rumus dasar yang tidak berasal dari kotak
\newcommand{\pakaidasar}[2]{\baris{Rumus}{{\small\sffamily\color{black!60}[#1]}\enspace$\displaystyle #2$}}
\newcommand{\susun}[2][]{\baris{Susun}{$\displaystyle #2$\if\relax\detokenize{#1}\relax\else
  \enspace{\small\color{black!60}(#1)}\fi}}
\newcommand{\angka}[1]{\baris{Angka}{$\displaystyle #1$}}
\newcommand{\hasil}[1]{\baris{Hasil}{\colorbox{biru3!10}{$\displaystyle #1$}}}
\newcommand{\hasilt}[1]{\baris{Hasil}{#1}}
\newcommand{\ceks}[1]{\baris{Cek}{#1}}

% Peta warna bab
\newcommand{\swatch}[1]{\textcolor{#1}{\rule{0.9em}{0.9em}}}
\newcommand{\swatches}[1]{\foreach \c in {#1}{\swatch{\c}\,}}
\newtcolorbox{petawarna}[1]{enhanced,colback=white,colframe=black!25,boxrule=0.4pt,arc=1pt,
  title={\faPalette\enspace Peta warna #1 \ (berlaku hanya di bab ini)},
  colbacktitle=black!6,coltitle=black,fonttitle=\bfseries\small,
  left=4pt,right=4pt,top=2pt,bottom=2pt,fontupper=\small}
\newenvironment{pw}{\begin{tabular}{@{}l>{\raggedright\arraybackslash}p{0.32\linewidth}>{\raggedright\arraybackslash}p{0.47\linewidth}@{}}}{\end{tabular}}
\newcommand{\pwrow}[3]{\swatches{#1} & \textbf{#2} & #3 \\}

% Kamus variabel
\newenvironment{kamus}{\par\medskip\textbf{Kamus variabel}\par\nopagebreak\small
  \setlength{\LTpre}{3pt}\setlength{\LTpost}{6pt}
  \begin{longtable}{@{}>{\centering\arraybackslash}p{2.0cm}>{\raggedright\arraybackslash}p{\dimexpr\linewidth-4.8cm-4\tabcolsep\relax}>{\raggedright\arraybackslash}p{2.8cm}@{}}
  \toprule Simbol & Nama dan arti & Satuan\\ \midrule \endhead
  \bottomrule \endfoot}{\end{longtable}\normalsize}

% Bab baru selalu di halaman baru
\newcommand{\bab}[1]{\clearpage\section{#1}}

% Panah fasor
\tikzset{fasor/.style={-{Stealth[length=2.6mm]},line width=1.1pt},
         bantu/.style={densely dashed,black!45}}

% Tata letak seksi
\usepackage{titlesec}
\titleformat{\section}{\Large\bfseries}{\thesection}{0.8em}{}
\titleformat{\subsection}{\large\bfseries}{\thesubsection}{0.7em}{}
\titleformat{\subsubsection}{\normalsize\bfseries}{\thesubsubsection}{0.6em}{}
\titlespacing*{\section}{0pt}{4pt}{10pt}

% ------------------------------------------------------------
%  Tambahan (dari TILI, dipakai juga di Proteksi)
% ------------------------------------------------------------
\usepackage{graphicx}
\graphicspath{{img/}}
\usepackage{needspace}
\usetikzlibrary{fit,backgrounds,shapes.misc,shapes.symbols}
\newcommand{\fig}[1]{\input{fig/#1}}
% beberapa foto sejajar: \fotos{tinggi}{file1/ket1,file2/ket2,...}
\newcommand{\fotos}[2]{\par\begin{center}\foreach \f/\k in {#2}{\begin{tabular}[b]{@{}c@{}}\includegraphics[height=#1]{\f}\\[-1pt]
  {\footnotesize\color{black!60}\emph{\k}}\end{tabular}\hspace{6mm}}\end{center}\par}
% gambar di tengah, diperkecil otomatis bila lebih lebar dari baris
\newsavebox{\figbox}
\newcommand{\figc}[1]{\par\begin{center}\sbox{\figbox}{\input{fig/#1}}\ifdim\wd\figbox>\linewidth
  \resizebox{\linewidth}{!}{\usebox{\figbox}}\else\usebox{\figbox}\fi\end{center}\par}
% foto slide: \foto[lebar]{file}{keterangan}
\newcommand{\foto}[3][0.7\linewidth]{\par\begin{center}\begin{minipage}{\linewidth}\centering
  \includegraphics[width=#1,height=5.6cm,keepaspectratio]{#2}\par\vspace{1pt}
  {\footnotesize\color{black!60}\emph{#3}}\end{minipage}\end{center}\par}
% gaya node diagram
\newcommand{\nohyph}{\hyphenpenalty=10000\exhyphenpenalty=10000\relax}
\tikzset{
  blok/.style={draw=black!60,rounded corners=2pt,align=center,font=\footnotesize,minimum height=7mm,inner sep=3pt,fill=black!3,execute at begin node=\nohyph},
  ar/.style={-{Stealth[length=2.2mm]},black!65,thick},
  bus/.style={line width=2.2pt},
  lbl/.style={font=\scriptsize,align=center,execute at begin node=\nohyph},
}
% simbol kecil untuk diagram satu garis
\newcommand{\cbsym}[2]{\draw[thick,fill=white] (#1) ++(-0.13,-0.13) rectangle ++(0.26,0.26); \node[font=\tiny,right] at ($(#1)+(0.15,0)$) {#2};}
\newcommand{\trafosym}[1]{\draw[thick] ($(#1)+(0,0.17)$) circle (0.2); \draw[thick] ($(#1)+(0,-0.17)$) circle (0.2);}

% ------------------------------------------------------------
%  Tambahan khusus Proteksi
% ------------------------------------------------------------
\usepgfplotslibrary{groupplots}
% simbol sekering kecil (zig-zag) di garis horizontal: \fusesym{titik}
\newcommand{\fusesym}[1]{\draw[thick] ($(#1)+(-0.22,0)$) .. controls ($(#1)+(-0.11,0.32)$) and ($(#1)+(-0.02,0.32)$) .. ($(#1)$) .. controls ($(#1)+(0.02,-0.32)$) and ($(#1)+(0.11,-0.32)$) .. ($(#1)+(0.22,0)$);}
% sekering pada garis vertikal
\newcommand{\fusesymv}[1]{\draw[thick] ($(#1)+(0,0.22)$) .. controls ($(#1)+(0.32,0.11)$) and ($(#1)+(0.32,0.02)$) .. ($(#1)$) .. controls ($(#1)+(-0.32,-0.02)$) and ($(#1)+(-0.32,-0.11)$) .. ($(#1)+(0,-0.22)$);}
% gelembung arus gangguan max/min
\newcommand{\gel}[3]{\node[draw,ellipse,inner sep=0.6pt,font=\tiny,align=center] at (#1) {#2\\[-1pt]\rule{4.5mm}{0.3pt}\\[-1pt]#3};}
% Pelengkap di luar slide: latar abu-abu, garis kiri tipis
\newtcolorbox{luar}{enhanced,breakable,colback=black!3,colframe=coklat4,
  boxrule=0pt,leftrule=1.5pt,arc=0pt,outer arc=0pt,fontupper=\small,
  title={\faBookOpen\enspace Pelengkap (di luar slide)},colbacktitle=black!3,coltitle=coklat1,
  fonttitle=\bfseries\small,left=7pt,right=7pt,top=3pt,bottom=3pt,before skip=6pt,after skip=6pt}
% kilat gangguan
\newcommand{\kilat}[1]{\draw[merah1,thick,-{Stealth[length=1.6mm]}] ($(#1)+(0.08,0.30)$)--($(#1)+(-0.06,0.05)$)--($(#1)+(0.08,0.02)$)--($(#1)+(-0.05,-0.26)$);}
